\scanpage{226}
(20) On a \unsure{équivalence|équivalences}
\[
\begin{aligned}
&f\text{ est une \uline{homothétie} pour }\gamma_\varphi
 &&\Longleftrightarrow\ \exists\lambda\in A,\text{ avec }f=\lambda_S,\\
&\text{i.e. }\forall(s,t)\in S\times S,\quad\{s,t\}\in A_\varphi\Longleftrightarrow
 \{f^{-1}(s),f(t)\}\in A_\varphi &&\text{et}\\
&\text{i.e. }\forall t\in S,\quad f(\mathrm{Et}_{A_\varphi}(t))=\mathrm{Et}_{A_\varphi}(f^{-1}(s))
 &&\lambda_{S^\vee}\varphi=\varphi\lambda_S^{-1}
\end{aligned}
\]
On trouve ainsi
\[
 \begin{aligned}
 &\text{Correspondance} &\text{homothéties} &&\text{Automorphismes }\lambda\in A\\
 &\text{biunivoque} &\text{du graphe} &\simeq&\text{du prégraphe }\Gamma,\text{ tels que}\\
 &\text{entre} &\gamma_\varphi=(S,A_\varphi) &&\lambda_{S^\vee}\varphi=\varphi\lambda_S^{-1}.
 \end{aligned}\tag*{(21)}
\]

On dit qu'une jauge $\varphi$ \deleted{\unsure{est une}} (sans supposer \unsure{néc.t|ici} $\Gamma$
séparé en $S^\vee$ ni $\varphi$ surjective) est compatible avec
l'autom $\lambda$, ou est une $\lambda$-jauge, si elle satisfait
\[
 \lambda_{S^\vee}\varphi=\varphi\lambda_S^{-1}\qquad\text{i.e. }\lambda_{S^\vee}\varphi\lambda_S=\varphi
 \quad\deleted{\unsure{\varphi t}}\tag*{(22)}
\]
% slanted
\margin{\uline{NB} Si $\varphi$ bijectif, la condition \unsure{(23)|(22)} s'écrit aussi
$(22')\ \varphi(\lambda)=\lambda_{S^\vee}^{-1}$ ou
$(22'')\ \varphi^{-1}(\lambda_{S^\vee})=\lambda_S^{-1}$}
(de $=$, si $Z\subset A$, on définit les $Z$-jauges
comme les jauges satisfaisant (22) $\forall\lambda\in Z$).
Ainsi, Si $\varphi$ est une $\lambda$-\unsure{gauche|jauche|jauge}, alors \deleted{\unsure{\varphi}} $\lambda_S$ est
une homothétie pour $\gamma_\varphi$, et inversement dans
le cas où $\Gamma$ est séparé en $S^\vee$.

\deleted{Supposons} On a
\[
\begin{aligned}
&\text{Si }\varphi\text{ est une }\deleted{\text{jau}}\{\lambda,\mu\}\text{-jauge }\inserted{\text{bijective}},
 \text{ où }\lambda,\mu\in A,\\
&\text{alors }\varphi\text{ est une }\lambda\mu\text{-jauge sss }\lambda\mu=\mu\lambda
\end{aligned}\tag*{(23)}
\]
ce qui implique que \deleted{pour} \inserted{deux homothéties d'}un graphe séparé,
\deleted{\unsure{deux}} \inserted{leur composé} \inserted{est une} homothétie sss elles commutent.
Si $Z$ est un sss-groupe de $A$, tel qu'il existe une $Z$-jauge
bijective, alors $Z$ doit être commutatif.

Donnons nous maintenant \deleted{\unsure{un}}
\[
 Z\subset A=\mathrm{Aut}(\Gamma)\qquad\text{un sous-groupe \uline{commutatif}}
 \tag*{(24)}
\]
et considérons \deleted{une $Z$-jauge $\varphi_0$.}
% archive sheet number: 225

